\chapter{Literature Review}

This literature review surveys peer-reviewed publications spanning synthetic data generation, tabular data modeling, generative adversarial networks, variational autoencoders, diffusion models, and privacy-preserving mechanisms. A systematic analysis of 151 research papers across 10 categories was conducted to identify the state of the art, evaluate existing limitations, and pinpoint the critical research gap addressed by our proposed approach. Table~\ref{tab:lit_categories} summarizes the scope of this review.

This review spans generative modeling, privacy, evaluation, and financial applications because synthetic tabular data sits at the intersection of all four. Understanding RE-TabSyn's position requires knowing what came before, where those approaches fell short, and what gap remains.

\begin{table}[htbp]
\centering
\caption{Summary of Literature Review Scope: 151 Papers Across 10 Categories}
\label{tab:lit_categories}
\small
\begin{tabularx}{\textwidth}{|l|c|X|}
\hline
\textbf{Category} & \textbf{Count} & \textbf{Key Papers} \\
\hline
A. Diffusion-Based Models & 21 & TabSyn, TabDDPM, CoDi, FinDiff, TabDiff \\
\hline
B. Privacy-Focused Models & 21 & DP-SGD, DP-CTGAN, PrivSyn, PATE-GAN \\
\hline
C. GAN-Based Models & 19 & CTGAN, CTAB-GAN+, TabFairGAN, WGAN \\
\hline
D. Transformer/LLM Models & 8 & REaLTabFormer, TabLLM, TabNet \\
\hline
E. Evaluation \& Benchmarking & 21 & SynthEval, SDMetrics, SynthCity \\
\hline
F. Privacy Attacks \& Defense & 17 & MIA Studies, TAMIS, Attribute Inference \\
\hline
G. Domain Applications & 18 & EHR-Safe, FinSyn, RareGraph-Synth \\
\hline
H. Theoretical Foundations & 16 & VAE, Copulas, Optimal Transport \\
\hline
I. Recent Advances & 8 & Foundation Models, Neural ODEs \\
\hline
J. Uncategorized & 2 & FedAvg, CLIP \\
\hline
\textbf{Total} & \textbf{151} & \\
\hline
\end{tabularx}
\end{table}

\section{Foundational Works in Generative Models}

The field of synthetic data generation is built upon three foundational paradigms in deep generative modeling, each offering distinct advantages and limitations. This section reviews the theoretical underpinnings and practical characteristics of Generative Adversarial Networks, Variational Autoencoders, and Denoising Diffusion Probabilistic Models -- the three pillars upon which modern tabular synthesis methods, including RE-TabSyn, are constructed.

In brief: a generative model is a system that learns what real data looks like by studying many examples, and then produces new data that resembles those examples. The challenge is that data is not simple. A table of financial records has dozens of columns, each with its own distribution, and they all relate to each other in complex ways. The three families of models below each take a different approach to solving this problem.

\subsection{Generative adversarial networks}

Goodfellow et al.\ \cite{goodfellow2014gan} introduced Generative Adversarial Networks, i.e., GANs, in 2014, establishing a framework wherein two neural networks, a generator and a discriminator, compete in a minimax game. The generator learns to produce samples indistinguishable from real data, while the discriminator learns to differentiate between real and generated samples. This adversarial training paradigm revolutionized generative modeling across domains, spawning thousands of follow-up works. At Nash equilibrium, the generator recovers the true data distribution, a theoretical guarantee that motivated widespread adoption.

A useful way to picture this: imagine a counterfeiter trying to produce fake currency and a bank inspector learning to spot the fakes. At the beginning, the counterfeiter's notes are obviously wrong, and the inspector catches them easily. Over thousands of rounds, the counterfeiter gets better at mimicking real currency, while the inspector gets sharper at detection. Eventually, the counterfeiter's notes are so convincing that even the inspector cannot tell them apart. That is, loosely, how a GAN works. The ``generator'' is the counterfeiter; the ``discriminator'' is the inspector. When training works well, the generator has learned to produce samples that match the real data distribution. In this research, instead of currency, the model is learning to produce realistic rows of financial data.

The Nash equilibrium mentioned above is a concept from game theory. It is the point at which neither player can improve by changing their strategy alone. For GANs, it represents the theoretical state where the generator perfectly replicates the data distribution. In practice, actually reaching this equilibrium is difficult, and this difficulty gives rise to the training problems described below.

However, GANs suffer from three well-documented pathologies:
\begin{enumerate}
    \item \textbf{Training Instability:} The minimax objective creates oscillatory dynamics where the generator and discriminator may fail to converge, requiring careful hyperparameter tuning. Small changes in learning rate or batch size can cause the entire training run to collapse.

    \item \textbf{Mode Collapse:} The generator may forget underrepresented modes of the data distribution, a critical failure for rare event generation. A GAN on a dataset with 5\% fraud can collapse to generating only the majority pattern, ignoring fraud entirely.

    \item \textbf{Convergence Diagnostics:} GAN training lacks a single meaningful loss metric. The training loss oscillates, and there is no obvious signal that the model has converged.
\end{enumerate}

Subsequent works addressed these issues: Wasserstein GAN \cite{arjovsky2017wgan} introduced the Earth Mover's distance as a more stable training objective, providing a gradient signal that is meaningful even when the generator and discriminator are far apart in distribution. Gulrajani et al.\ \cite{gulrajani2017wgangp} improved WGAN training further with a gradient penalty term that enforces a Lipschitz constraint on the discriminator, avoiding the weight-clipping instabilities of the original WGAN. Salimans et al.\ \cite{salimans2016improved} proposed several training techniques: feature matching (training the generator to match intermediate representations in the discriminator rather than its final output), minibatch discrimination (allowing the discriminator to look at multiple samples simultaneously to detect mode collapse), and historical averaging (penalizing the model for deviating too much from the average of its past parameters). CausalGAN \cite{kocaoglu2018causalgan} introduced causal reasoning into adversarial training, allowing generation conditioned on causal interventions rather than mere correlations. Despite these advances, mode collapse on minority classes remains a fundamental limitation for imbalanced data generation, and it is one of the key reasons this research turns to diffusion models instead.

\subsection{Variational autoencoders}

The Variational Autoencoder, i.e., VAE, \cite{kingma2014vae}, introduced by Kingma and Welling in 2014, takes a different approach to generative modeling. Rather than pitting two networks against each other, a VAE learns to compress data into a compact representation and then reconstruct it. The ``variational'' part refers to the fact that this compression is probabilistic: instead of mapping each data point to a single point in a compressed space, the VAE maps it to a probability distribution (specifically, a Gaussian distribution parameterized by a mean and a variance). This smooth, probabilistic compression is what makes generation possible.

To illustrate: imagine storing a description of every person in a country using just a few numbers. You might encode each person using ten numbers that capture age, height, occupation, income, and so on. A VAE learns to find the best set of numbers (the ``latent representation'') such that, given those numbers, you can reconstruct the original person's description as accurately as possible. The key insight is that nearby points in this compressed space correspond to similar people, so you can generate a new, realistic person by picking a random point in the compressed space and ``decoding'' it.

Formally, the VAE combines variational inference with deep learning, enabling probabilistic latent variable modeling. By maximizing the Evidence Lower Bound, i.e., ELBO, VAEs learn smooth latent representations suitable for generation. The ELBO has two parts: a reconstruction term that encourages the model to faithfully reproduce the input, and a KL divergence term that regularizes the latent space to be close to a standard Gaussian distribution. The reparameterization trick enables end-to-end gradient-based optimization, making VAEs both theoretically principled and practically trainable. Without the reparameterization trick, the sampling step in the VAE would be non-differentiable, making gradient descent impossible. The trick reformulates the sampling as a deterministic function of the mean, the standard deviation, and a separate random noise variable, allowing gradients to flow through the mean and variance parameters while the randomness comes from outside the computation graph.

VAEs offer several advantages over GANs. Training is stable because the ELBO objective is well-defined and monotonically improvable, avoiding oscillatory adversarial dynamics. The KL divergence regularization produces a smooth, continuous latent geometry where interpolation between points yields semantically meaningful samples. VAEs also provide an approximate likelihood, enabling model comparison and anomaly detection — useful in finance for flagging transactions whose latent representation is far from the learned distribution.

However, VAEs often suffer from blurry reconstructions due to the Gaussian decoder assumption (which averages over modes) and posterior collapse in complex datasets (where the model ignores the latent code entirely and relies solely on the decoder). Posterior collapse is particularly problematic: the model stops using the compressed representation and learns to reconstruct everything from scratch in the decoder, defeating the purpose of the latent space. Kim et al.\ \cite{kim2023tabvae} proposed Tab-VAE, a specialized VAE for tabular data that addresses some of these limitations through mixed-type decoding heads that treat numerical and categorical features differently, though it still lacks class conditioning capabilities.

RE-TabSyn uses a VAE as its first component, but with specific design choices (including a $\beta$-VAE formulation with $\beta = 0.1$) that prevent posterior collapse while maintaining a smooth latent geometry suitable for downstream diffusion. This is discussed in detail in Chapter~2.

\subsection{Denoising diffusion probabilistic models}

Ho et al.\ \cite{ho2020ddpm} presented Denoising Diffusion Probabilistic Models, i.e., DDPMs, in 2020, and they reshaped generative modeling almost overnight. The central idea is both simple and powerful: start with real data, gradually add random noise until the data is completely destroyed, and then train a model to reverse that process step by step.

Think of it this way. Take a clear photograph and start adding random static to it, a little at a time, until after a thousand steps the image is just pure noise with no trace of the original picture. Now train a neural network to do the reverse: given a slightly noisy image, predict what it looked like one step earlier, with slightly less noise. If this network works, you can start from pure noise and walk backwards step by step, producing a realistic image. Critically, you can influence what image gets generated by conditioning each denoising step on some additional information, such as a text description or a class label.

DDPMs achieved a breakthrough in image generation, surpassing GANs in both fidelity (measured by FID score) and diversity (measured by recall metrics). The fact that diffusion models do not suffer from mode collapse, a direct consequence of the denoising objective that must cover all modes to accurately reverse noise, made them particularly attractive for tasks involving rare or minority classes.

Diffusion models offer three key advantages relevant to this work:
\begin{enumerate}
    \item \textbf{Stable Training:} The mean squared error denoising objective provides consistent gradient signals without adversarial instabilities.

    \item \textbf{Mode Coverage:} The denoising objective must cover all modes of the data distribution, including rare classes, to keep denoising error low. Unlike GANs, diffusion models do not suffer from mode collapse.

    \item \textbf{Flexible Conditioning:} Conditioning information can be incorporated at each denoising step, enabling fine-grained control over generated outputs — the property RE-TabSyn exploits through CFG.
\end{enumerate}

Song et al.\ \cite{song2021score} provided a unified theoretical framework connecting DDPMs with score-based models through stochastic differential equations, i.e., SDEs, establishing that both approaches learn the score function $\nabla_\mathbf{x} \log p(\mathbf{x})$ of the data distribution. The score function points in the direction of increasing data density. Learning it means learning where data is likely to come from, which is exactly what you need for generation. The foundational thermodynamic perspective was first articulated by Sohl-Dickstein et al.\ \cite{sohldickstein2015deep}, and Dhariwal and Nichol \cite{dhariwal2021diffbeatgan} demonstrated convincingly that diffusion models surpass GANs on image synthesis benchmarks, putting to rest any remaining doubts about their generational quality.

Rombach et al.\ \cite{rombach2022ldm} introduced Latent Diffusion Models, i.e., LDMs, demonstrating that performing diffusion in a compressed latent space reduces computational cost significantly while maintaining generation quality. The core insight: instead of running the diffusion process directly on high-dimensional data (like raw pixel values or raw tabular features), first compress the data into a smaller latent representation using a VAE, then run diffusion in that compressed space. This decouples the expensive perceptual compression step from the generative modeling step. This insight, operating diffusion in a learned latent space rather than the raw data space, is central to RE-TabSyn's design. It is precisely why RE-TabSyn succeeds where direct diffusion approaches like TabDDPM fail, a point that becomes clearer in the next section.

Table~\ref{tab:gen_model_comparison} summarises the key properties of the three families and their practical consequences for tabular synthesis.

\begin{table}[htbp]
\centering
\caption{Comparison of Generative Model Families for Tabular Synthesis}
\label{tab:gen_model_comparison}
\small
\begin{tabularx}{\textwidth}{|l|C|C|C|}
\hline
\textbf{Property} & \textbf{GANs: CTGAN} & \textbf{VAEs: TVAE} & \textbf{Diffusion: TabSyn, RE-TabSyn} \\
\hline
Training stability & Low (adversarial) & High (ELBO) & High (MSE denoising) \\
\hline
Mode coverage & Poor (mode collapse) & Moderate & Excellent \\
\hline
Sample diversity & Low to moderate & High & High \\
\hline
Reconstruction sharpness & High & Blurry (Gaussian avg.) & High \\
\hline
Latent space quality & None & Smooth Gaussian & Smooth Gaussian (via VAE) \\
\hline
Evaluation metric & FID / KS & Reconstruction loss / KS & KS / TSTR AUC \\
\hline
Training time (relative) & 1$\times$ & 0.5$\times$ & 2--4$\times$ \\
\hline
Class conditioning support & Conditional GAN & None (standard) & CFG (RE-TabSyn only) \\
\hline
Minority class performance & Poor (mode collapse) & Moderate & Strong (RE-TabSyn) \\
\hline
Formal DP compatibility & Difficult (mode collapse worsens) & Moderate & Moderate to good \\
\hline
\end{tabularx}
\end{table}

The table makes clear why the diffusion paradigm is the right foundation for RE-TabSyn. GANs have excellent sample quality when they work, but the mode collapse problem specifically targeting minority classes is fatal for the financial imbalanced data problem. VAEs are stable but produce blurry reconstructions and have no class conditioning. Diffusion models combine high sample quality, good mode coverage, and a natural conditioning mechanism via CFG that enables RE-TabSyn's core contribution.

\section{Synthetic Tabular Data Generation}

Building on these foundational generative paradigms, a growing body of work has specifically targeted the challenge of synthesizing realistic tabular data. Tabular data presents challenges that simply do not exist for images or text. A single row in a financial table might have 30 columns: some are continuous numbers (income, transaction amount), some are categories (loan type, employment status), some have missing values, and all of them are related to each other in complex ways (high income tends to correlate with low default rates, married people with dependents tend to have specific spending patterns, and so on). None of the image-generation techniques transfer directly to this setting.

This section traces the major tabular synthesis methods in chronological order, showing how each generation of models addressed the limitations of its predecessors, and where each one ultimately fell short.

\subsection{CTGAN and TVAE}

CTGAN \cite{xu2019ctgan}, proposed by Xu et al.\ in 2019, was the first serious attempt to adapt the GAN framework specifically for tabular data. It addresses two core challenges of tabular synthesis. First, financial numerical data often has complex distributions that are nothing like a simple bell curve. Income data, for instance, tends to have multiple ``peaks'' corresponding to different income brackets, with long tails for very high earners. Standard normalization methods fail to capture this shape. CTGAN handles this through mode-specific normalization for numerical columns, using variational Gaussian mixture models to explicitly model and transform these multimodal distributions before feeding them into the GAN.

Second, tabular data often has imbalanced categorical distributions. In a dataset of loan applications, the vast majority might be for mortgages, with only a small fraction for personal loans. Without special handling, the GAN would learn to generate mortgage applications almost exclusively. CTGAN addresses this through a conditional generator with training-by-sampling: a conditional vector specifies which categorical value to generate, and the training procedure oversamples rare categories. This means the discriminator sees a balanced distribution of categories during training, giving the generator an equal incentive to learn all of them.

CTGAN became the de facto baseline for tabular synthesis, implemented in the widely-used Synthetic Data Vault, i.e., SDV, library and reproduced in hundreds of subsequent papers.

Introduced alongside CTGAN, TVAE applies variational autoencoders to tabular synthesis using the same preprocessing pipeline but ELBO optimization instead of adversarial training. TVAE offers more stable training at the cost of slightly lower sample quality, making it a reliable fallback when CTGAN training proves unstable. Conditional Wasserstein GAN \cite{engelmann2021cwgan} explored oversampling for imbalanced learning, and TabFairGAN \cite{rajabi2022tabfairgan} addressed fairness constraints by including fairness terms in the GAN objective. VAE-GAN hybrid architectures \cite{camino2020vaegan, pereira2024synthetic} combined strengths of both paradigms, while Tab-VAE \cite{kim2023tabvae} specialized the VAE framework for tabular data. DP-CTGAN \cite{torfi2022dpctgan} and Fed-DPSDG-WGAN \cite{ma2023feddpsdg} extended GAN-based synthesis with privacy and federated constraints. Esteban et al.\ \cite{esteban2017rcgan} applied conditional GANs to medical time series, and hybrid deep learning approaches \cite{mishra2022hybrid, yoon2020gain} addressed domain-specific challenges. Beyond the Norm \cite{nidhi2024beyond} surveyed rare event generation comprehensively. Both CTGAN and TVAE share the same limitation: they faithfully replicate the training class distribution without any mechanism for adjustment. The model will generate approximately the same fraction of fraud cases as exists in the training data, because nothing in the training objective tells it to do anything different.

\subsection{CTAB-GAN+}

Zhao et al.\ \cite{zhao2022ctabgan} enhanced CTGAN with auxiliary classifiers and information-maximizing discriminators in CTAB-GAN+. The auxiliary classifier provides an additional gradient signal that helps the generator produce samples from the correct class, while the information-maximizing discriminator encourages the generator to fully utilize its input noise, reducing mode collapse. This model improved statistical similarity and downstream utility, particularly for datasets with complex feature interactions such as those found in healthcare records and financial risk assessments. Several ablation studies confirmed that both components contribute independently to the improvement.

However, CTAB-GAN+ remained bound to the training distribution. The auxiliary classifier helps each class be represented more accurately, but it does not allow the user to request more minority samples than exist in the training data. The ratio of classes in the generated data still mirrors the ratio in the real data. It introduced additional training complexity without addressing the class imbalance control problem that motivates RE-TabSyn.

\subsection{TabDDPM}

TabDDPM \cite{kotelnikov2023tabddpm}, published at ICML 2023 by Kotelnikov et al., was the first systematic application of diffusion models to tabular data. The paper matters both for its methodology and for making explicit the failure modes of applying diffusion naively to heterogeneous data.

TabDDPM's approach: treat numerical features with Gaussian diffusion (the standard continuous noise process) and categorical features with multinomial diffusion (a discrete version that treats category assignments as probability vectors over $K$ classes). While it demonstrated competitive performance on purely numerical datasets, TabDDPM exhibited catastrophic failure on mixed-type financial data.

Why does it fail? The problem lies in how categorical features are represented. One-hot encoding, the standard method, turns a category like ``self-employed'' into a vector like $[0, 0, 1, 0, 0]$. When you add Gaussian noise to this vector, you get something like $[0.03, -0.12, 0.91, 0.07, -0.02]$, which does not correspond to any valid category. The model must learn to denoise vectors that never appear in the real data, creating sparse, discontinuous gradients that make training very difficult. This is analogous to trying to learn to draw faces by starting from a blurry mess that could not possibly be a face at any intermediate step.

Our experiments confirm this failure: TabDDPM achieves an average KS statistic of 0.770 across six financial datasets, far worse than even simple GAN baselines. A KS statistic of 0.770 means the synthetic distributions are almost completely different from the real distributions; the model is generating random noise. This failure validates the latent-space approach that both TabSyn and RE-TabSyn adopt.

\subsection{TabSyn}

TabSyn \cite{zhang2024tabsyn}, published at ICLR 2024 by Zhang et al., achieves the highest reported fidelity for tabular data synthesis. It directly addresses TabDDPM's categorical encoding problem through a two-stage design:
\begin{enumerate}
    \item A VAE encodes the mixed-type tabular data (numerical and categorical columns together) into a smooth, continuous latent representation where all features are represented as real-valued vectors. The discreteness problem disappears: instead of working with raw category labels, the diffusion model works with the VAE's learned continuous encodings.
    \item A Transformer-based diffusion model operates in this smooth latent space, capturing complex inter-column dependencies through self-attention. The Transformer architecture, originally developed for natural language processing, is particularly well-suited here because tabular columns, like words in a sentence, have complex relationships with each other that are worth modeling explicitly.
\end{enumerate}

TabSyn achieves exceptional fidelity with an average KS statistic of 0.109 in our experiments, establishing a new benchmark for distributional similarity. On most datasets, the synthetic and real distributions are almost indistinguishable statistically. For a practitioner, this means that TabSyn synthetic data is nearly a perfect stand-in for real data when it comes to training machine learning models.

However, TabSyn has no mechanism for controlling class distribution whatsoever. It mirrors the training data's imbalance exactly. If a dataset has 5\% minority class, TabSyn generates approximately 5\% minority samples. For a fraud detection application where the goal is to train a better model on balanced data, TabSyn provides no benefit over using the real data directly. The core limitation is architectural: nothing in the training objective tells the model to weight any class differently.

\subsection{Other diffusion-based approaches}

The rapid success of TabSyn motivated several extensions addressing different aspects of the tabular synthesis problem. TabDiff \cite{shi2024tabdiff} proposes a multi-modal diffusion framework that jointly models numerical and categorical features without requiring a VAE encoding step, using a unified mixed-type noise process. FinDiff \cite{sattarov2023findiff} applies diffusion specifically to financial tabular data, incorporating domain-specific preprocessing and evaluation protocols adapted to financial data characteristics. RelDDPM \cite{liu2024relddpm} extends diffusion for relational multi-table synthesis, where multiple tables with foreign-key relationships must be synthesized simultaneously while preserving referential integrity.

CoDi \cite{choi2023codi} introduces co-evolving contrastive diffusion for mixed-type synthesis, using contrastive learning to align numerical and categorical representations during the denoising process. Zheng et al.\ \cite{zheng2024controllable} explore controllability through diffusion conditioning, though their approach requires auxiliary classifiers and cannot achieve the post-hoc ratio control that RE-TabSyn's CFG mechanism provides. Additional works address federated settings \cite{noble2024dpfed, sattarov2024fedtabdiff, yang2024feddiff, shankar2024silofuse}, domain-specific constraints \cite{yuan2024domain}, missing value imputation \cite{ouyang2023missdiff, jolicoeurmartineau2024diffusion}, and surveys of the tabular diffusion domain \cite{mueller2024diffusion}. Difformer \cite{wu2023difformer} and conditional score-based approaches \cite{baldassari2024conditional} further expand the theoretical foundations.

While each of these works makes valuable contributions to the field, none addresses the problem of controllable class-conditional generation, the ability to specify exactly what fraction of the generated data should belong to the minority class.

\section{Controllable and Conditional Generation}

A key limitation shared by all methods reviewed above is the inability to control the class distribution of generated samples. This might seem like a minor technical detail, but for financial applications it is the central problem. A fraud detection team that uses TabSyn to generate a million synthetic transactions will still end up with only 5,000 to 50,000 fraudulent ones, the same 0.5-5\% prevalence as in their real data. That does not help them train a better model. What they need is a way to generate, say, 400,000 fraudulent and 600,000 legitimate transactions, preserving the realistic statistical properties of fraud while greatly increasing its representation.

This section examines existing approaches to conditional and controllable generation, and explains why none of them fully solve this problem for tabular data.

\subsection{Classifier-free diffusion guidance}

Ho and Salimans \cite{ho2022cfg} introduced Classifier-Free Guidance, i.e., CFG, in 2022, and it quickly became the standard mechanism for controlled generation in image and text domains.

The key idea is elegant. Instead of training a separate classifier to guide generation (which requires extra training and introduces complex dependencies), CFG trains a single model that can operate either with or without class conditioning. During training, the class label is randomly replaced with a special ``null token'' with some probability (typically 10--20\%). The model therefore sees two types of training examples: ones where it knows which class to generate, and ones where it has no class information at all. It learns both a conditional distribution (``what does a fraudulent transaction look like?'') and an unconditional distribution (``what does a transaction look like in general?'').

At generation time, both distributions are queried simultaneously. The conditional prediction points in the direction of the desired class; the unconditional prediction represents the baseline. CFG amplifies the difference between the two, steering the generation more strongly toward the desired class than standard conditioning would:

\begin{equation}
\tilde{\epsilon} = \epsilon_{uncond} + w \cdot (\epsilon_{cond} - \epsilon_{uncond})
\end{equation}

The parameter $w$ is the guidance scale. At $w = 1$, you get standard conditional generation. At $w = 2$ (the default in RE-TabSyn), the model is pushed twice as far in the direction of the desired class, producing a strong class signal while still generating realistic-looking samples. This guidance scale provides an intuitive trade-off: higher $w$ produces outputs more aligned with the condition but potentially less diverse.

CFG quickly became the standard conditioning mechanism for text-to-image models (Stable Diffusion, DALL-E 2, Imagen) due to its simplicity and effectiveness. It requires no separate classifier training, no modifications to the base architecture, and can be applied at inference time without retraining. This last property is particularly important for RE-TabSyn: a practitioner can train the model once on the real imbalanced data and then generate datasets with any desired minority ratio at runtime.

Despite its demonstrated success in image generation, CFG had never been applied to tabular data synthesis prior to our work. RE-TabSyn is the first to bridge this gap, adapting CFG for structured data class control.

\subsection{RelDDPM and other conditional methods}

RelDDPM \cite{liu2024relddpm} extends diffusion models for controllable tabular synthesis through conditioning mechanisms that support inter-table relationships and target constraints. However, its conditioning architecture is complex, requires auxiliary classifiers, and does not provide the simple, continuous control that CFG offers. Specifically, RelDDPM's conditioning is designed for relational constraints (foreign keys, referential integrity) rather than for class ratio control.

Transformer-based approaches, including REaLTabFormer \cite{solatorio2023realtabformer}, TabLLM \cite{hegselmann2023tabllm}, TabNet \cite{arik2021tabnet}, SAINT \cite{somepalli2021saint}, TabTransformer \cite{huang2020tabtransformer}, and FLAIM \cite{liu2024flaim}, capture contextual dependencies through attention mechanisms, treating each column in a table row as analogous to a word in a sentence. Gorishniy et al.\ \cite{gorishniy2021revisiting} benchmarked deep learning models for tabular data, finding that well-tuned tree-based models (like XGBoost) often still outperform deep learning on tabular prediction tasks, though deep learning excels for generation. Flemings and Razaviyayn \cite{flemings2024private} addressed private prediction for large-scale text generation. Other conditional generation methods rely on conditional inputs to the generator but cannot modify the generation ratio after training: the conditioning governs what kind of sample is generated, not how many of each kind.

\section{Privacy-Preserving Synthetic Data}

Privacy concerns have driven substantial work on differentially private, i.e., DP, synthetic data generation. This is especially important in finance, where generating synthetic data that could be used to re-identify individuals would defeat the purpose entirely. This section explains the key concepts in privacy-preserving synthesis and situates them relative to RE-TabSyn.

Differential privacy, i.e., DP, is a mathematical framework for quantifying privacy guarantees. A mechanism is $(\epsilon, \delta)$-differentially private if running it on two datasets that differ by one person produces outputs whose distributions are indistinguishable (to within a factor of $e^\epsilon$ in probability, with a $\delta$ slack for unlikely events). Smaller $\epsilon$ means stronger privacy. The formal guarantee ensures that the participation of any single individual in the training data has a negligible effect on the synthetic data, making it mathematically impossible to conclusively identify whether a specific person's record was used in training.

DP-SGD \cite{abadi2016dpsgd}, introduced by Abadi et al.\ in 2016, formalized differentially private deep learning by clipping per-sample gradients and adding calibrated Gaussian noise during training. The clipping prevents any individual's data from dominating the gradient updates; the noise ensures indistinguishability. This provides formal $(\epsilon, \delta)$-differential privacy guarantees but significantly degrades model utility due to the noisy gradient updates. The trade-off is fundamental: more noise means better privacy but worse model quality.

PATE-GAN \cite{jordon2019pategan} combines the Private Aggregation of Teacher Ensembles, i.e., PATE, framework with GANs. Multiple ``teacher'' discriminators are trained on disjoint subsets of data. Their votes on whether a generated sample is real or fake are aggregated with Laplace noise, and only this noisy aggregate is used to train the student discriminator. This limits the information flow from any individual teacher's private data to the publicly trained generator.

DP-CTGAN \cite{torkzadehmahani2019dpctgan} integrates DP-SGD directly into CTGAN training, providing formal privacy guarantees. However, the noise injection compounds CTGAN's existing mode collapse problem: when gradient signals are noisy, the generator finds it even harder to learn rare minority patterns, making minority class generation even worse. Yale et al.\ \cite{yale2020synthetic} demonstrated that privacy-preserving generation significantly degrades utility for health data, a finding corroborated across financial datasets in multiple independent studies.

Stadler et al.\ \cite{stadler2022synthetic} raised important concerns in their ``Synthetic Data -- Anonymisation Groundhog Day'' paper, demonstrating that synthetic data without formal DP guarantees may still be vulnerable to membership inference attacks (where an adversary tries to determine whether a specific record was in the training data). Their analysis showed that statistical similarity alone is not a sufficient privacy guarantee: a model that perfectly replicates the training distribution is most exposed to inference attacks. Dwork and Roth \cite{dwork2014algorithmic} established the algorithmic foundations of differential privacy, while Dwork et al.\ \cite{dwork2006calibrating} introduced noise calibration techniques that set the standard for how much noise to add to achieve a given privacy level.

Additional privacy-focused works include data synthesis with missing data \cite{ge2024dpgen, park2018data}, differentially private diffusion models \cite{dockhorn2023dpddpm}, Gaussian process regression under DP \cite{smith2021dpreg}, foundation model APIs for DP synthesis \cite{ghalebikesabi2023dpsynth}, LLM-based DP tabular synthesis \cite{touati2024dptabllm}, text anonymization via DP-VAE \cite{weggenmann2022dpvae}, evaluation of DP in high-stakes domains \cite{ganev2023dpeval, hittmeir2020evaluating}, federated learning with matched averaging \cite{wang2020fedma}, tabular DP generation \cite{liu2024gentab}, the Opacus library \cite{yousefpour2021opacus}, practical DP in deep learning \cite{mcmahan2017practical}, privacy-preserving deep learning releases \cite{chen2020ppsynth}, private synthetic data for marginal queries \cite{liu2021privsynth}, PrivSyn \cite{zhang2021privsyn}, adversarial demographic attribute removal \cite{elazar2018adversarial}, and comprehensive surveys of DP tabular synthesis \cite{liu2024dptabsurvey}. This body of work motivates the inclusion of privacy metrics in our evaluation framework.

RE-TabSyn does not claim formal DP guarantees in its current form, but it is evaluated using the Distance to Closest Record, i.e., DCR, metric, which measures empirical privacy. This distinction matters: DCR tells you whether your synthetic data is memorizing training records in practice; DP tells you what an adversary can theoretically infer. Both perspectives are important, and future work integrating DP-SGD into RE-TabSyn's training is discussed in Chapter~7.

\section{Financial Domain Applications}

The financial sector has been an active adopter of synthetic data technology, driven by the unique combination of strict regulations and critical analytical needs that was discussed at length in Chapter~1. It is worth reviewing the specific published work in this area to understand what has been tried, what has worked, and where the remaining gaps are.

Assefa et al.\ \cite{assefa2020synthetic} provided an early survey of synthetic data opportunities and challenges in finance, identifying fraud detection, credit scoring, and stress testing as key applications. Their survey made clear that the bottleneck was not awareness of the technology but rather the quality gap: existing synthetic data methods could not reliably preserve the complex statistical properties that financial ML models depend on. Fiore et al.\ \cite{fiore2019hybrid} demonstrated that GAN-augmented training data improved credit card fraud detection accuracy by up to 10\%, establishing an early empirical precedent for synthetic data utility in imbalanced financial tasks. The 10\% improvement, while modest, was enough to justify continued investment in the research direction.

FinDiff \cite{sattarov2023findiff} applied diffusion models specifically to financial tabular data generation, achieving strong statistical similarity but without addressing the class control problem. Their work confirmed that diffusion models generalize well to financial data when the preprocessing is handled correctly, validating the approach taken by TabSyn and RE-TabSyn. Chen et al.\ \cite{chen2023finsyn} proposed FinSyn with constraint validation for financial data, incorporating domain-specific business rule checks (such as ensuring that generated loan amounts do not exceed simulated collateral values) into the GAN training loop. Their work highlighted that statistical fidelity alone is insufficient for financial applications: the synthetic data must also satisfy structural constraints that domain experts can verify.

Nidhi et al.\ \cite{nidhi2024rare} published a survey on synthetic data generation for rare event prediction, and their conclusion is worth quoting directly: ``the ability to control the generation of rare events remains an open challenge.'' This survey, published in 2024, provides independent confirmation that the research gap targeted by RE-TabSyn is real, open, and recognized by the community.

Domain-specific synthetic data applications span healthcare \cite{theodorou2023ehrsafe, yuan2023scoehr, li2024ehr, li2024ehrdiff, gonzales2023synthetic, nikolenko2021synthetic, goyal2024raregraph, dahmen2019geneval}, cybersecurity \cite{dasgupta2024review, min2024novel}, education \cite{luo2024fairsynedu}, energy systems \cite{ge2024contextual}, robotics \cite{mahler2017dexnet}, financial time series \cite{takahashi2019gansfinance, sattarov2024leveraging}, and the foundational Synthetic Data Vault framework \cite{patki2016synthvault}. Our own work \cite{shroff2026retabsyn}, accepted at the I2IT International Conference, directly addresses the open challenge identified by Nidhi et al.

\section{Evaluation Frameworks and Metrics}

Evaluating synthetic data quality is harder than it might seem. A naive approach might ask: does the synthetic data look similar to the real data? But ``looking similar'' can mean many different things. The distributions could match for each feature individually while pairwise correlations between features are completely wrong. Or the distributions and correlations could both match while a machine learning model trained on the synthetic data completely fails to generalize to real data. Or all of the above could be fine while some synthetic samples are exact copies of real records, creating a privacy risk.

The literature has converged on three pillars of evaluation, each capturing a different dimension of quality.

Statistical fidelity measures how closely synthetic data matches real data distributions. Key metrics include the Kolmogorov-Smirnov, i.e., KS, statistic for numerical columns, the chi-square test \cite{pearson1900criterion} for categorical columns, and the Jensen-Shannon Divergence, i.e., JSD, for overall distributional comparison. The KS statistic computes the maximum absolute difference between the cumulative distribution functions of the real and synthetic data for a given numerical column. A KS statistic of 0 means the distributions are identical; a value of 1 means they are completely non-overlapping. In practice, values below 0.20 indicate good fidelity, and values above 0.50 indicate distributional mismatch. Figueira and Vaz \cite{figueira2022survey} provide a comprehensive taxonomy of fidelity metrics.

Beyond marginal distributions, correlation preservation measures how well the synthetic data captures the relationships between features. Two features might each be individually well-distributed but their correlation in the synthetic data might differ from the real data. For example, in a credit risk dataset, income and credit score are positively correlated in reality; synthetic data that fails to capture this might generate high-income records with poor credit scores that have no real-world analogue.

Machine learning utility measures whether synthetic data can effectively replace real data for ML tasks. The Train-on-Synthetic, Test-on-Real, i.e., TSTR, protocol is the gold standard: train a model exclusively on synthetic data and evaluate it on held-out real data. If the synthetic data captures the relevant statistical structure, the model should perform comparably to one trained on real data. Metrics include AUC-ROC (a threshold-independent measure of classification ability), F1-score for the minority class (which captures both precision and recall for rare events), and accuracy for overall performance. The gap between ``Train on Real, Test on Real'', i.e., TRTR, performance and TSTR performance quantifies the utility ratio: how much of the real data's predictive value the synthetic data preserves. Jordon et al.\ \cite{jordon2022fake} established guidelines for effective utility evaluation, emphasizing that the TSTR protocol must be accompanied by minority-class metrics in imbalanced settings.

Privacy preservation measures resistance to re-identification attacks. Distance to Closest Record, i.e., DCR, quantifies how far each synthetic sample is from its nearest real neighbor in the feature space, with higher values indicating stronger empirical privacy. If the minimum DCR approaches zero, it suggests that some synthetic samples are nearly identical to real records, which is a privacy concern. Membership Inference Attack, i.e., MIA, resistance measures whether an adversary can determine if a specific record was in the training data, typically by training a classifier to distinguish between training members and non-members based on the model's behavior. Key works in privacy attacks include foundational MIA studies \cite{shokri2017mia, carlini2022mia, yeom2018privacy}, MIA against synthetic data \cite{vanbreugel2023mia, chen2020miahealth, hu2024empirical}, attribute inference attacks \cite{giomi2023linear}, unified privacy risk frameworks \cite{giomi2023unified}, model inversion attacks \cite{fredrikson2015model}, property inference \cite{ganju2018property}, re-identification estimation \cite{rocher2019estimating}, graph privacy \cite{wu2022linkteller}, de-anonymization \cite{narayanan2008robust}, TAMIS \cite{meeus2024tamis}, the MIDST challenge \cite{meeus2024midst}, and quantifying privacy risks for generative models \cite{stadler2023quantifying}. Additional evaluation works include universal metrics \cite{platzer2021universal}, privacy-utility assessment \cite{emam2022assessment}, DP benchmarks \cite{tao2021benchmark, liu2024benchmarking}, fairness studies \cite{ganev2024fair}, deep learning for tabular data surveys \cite{borisov2022deep}, mutual information estimation \cite{kraskov2004mutual}, SDV quality metrics \cite{patki2016sdv}, user needs analysis \cite{kolsch2024navigating}, scalability \cite{wang2024scaling}, survey data evaluation \cite{snoke2018synthetic}, attack-defense surveys \cite{bellovin2024synthtab}, SynthEval \cite{lautrup2024syntheval}, systematic assessment \cite{chundawat2022systematic}, the KS test \cite{massey1951ks}, and post-hoc inference \cite{wang2024unlocking}.

This three-pillar framework, fidelity, utility, and privacy, is the evaluation standard adopted in this research. The key innovation in our evaluation setup is the fourth dimension: minority class control, which no existing framework had previously needed to address because no existing method could control minority ratios. Chapter~4 describes how all four dimensions are measured in practice.

\section{Baseline Selection Criteria}

Any comparative study requires a deliberate choice of baselines. Including too few makes the comparison insufficient; including too many obscures the narrative. RE-TabSyn is compared against exactly four baselines: CTGAN, TVAE, TabDDPM, and TabSyn. Each was chosen for a specific reason, and together they span the space of relevant prior work in a principled way.

\subsection{CTGAN as GAN representative}

CTGAN \cite{xu2019ctgan} is the industry standard for tabular data synthesis. It has been cited over 2,000 times, is the default model in the Synthetic Data Vault, i.e., SDV, library that financial institutions use in practice, and has been reproduced independently across hundreds of publications. Any claimed improvement in tabular synthesis must beat CTGAN, or the comparison is incomplete. Including it also anchors the fidelity axis of the comparison: CTGAN achieves KS $\approx$ 0.153, which is a known reference point.

CTGAN represents the GAN family. If RE-TabSyn outperforms it, the result validates the diffusion-over-GAN design choice for this task.

\subsection{TVAE as VAE representative}

TVAE was introduced in the same paper as CTGAN \cite{xu2019ctgan}. It uses a VAE with the same preprocessing pipeline as CTGAN but replaces the adversarial objective with a variational one. Including TVAE serves two specific purposes:

First, it isolates the effect of the generative model family (VAE versus GAN) from the effect of preprocessing. Since CTGAN and TVAE share the same preprocessing, any performance difference between them reflects the difference between adversarial and variational training. This ablation is useful for understanding the relative contribution of each component.

Second, RE-TabSyn's first component is also a VAE. Comparing against TVAE tests whether RE-TabSyn's VAE-plus-diffusion design is strictly better than a VAE alone. The answer is yes -- TVAE achieves weaker fidelity and no class control -- but the comparison needs to be made explicitly.

\subsection{TabDDPM as diffusion failure case}

TabDDPM \cite{kotelnikov2023tabddpm} is included as a critical negative baseline. Its inclusion is not to show a close competitor but to demonstrate a failure mode that RE-TabSyn's architecture specifically addresses.

TabDDPM was the first paper to seriously apply diffusion models to tabular data. It was published at ICML 2023, a top-tier venue, and attracted considerable attention. If it had not been included in the comparison, a reader might reasonably ask: ``If diffusion works so well for images and text, why not just apply it directly to tables?'' TabDDPM answers that question empirically: direct application fails (KS = 0.770), and the reason is the one-hot categorical encoding mismatch that the VAE latent space solves.

Without TabDDPM, RE-TabSyn's architectural choice of latent-space diffusion would appear to be one option among many. With TabDDPM in the comparison, it is shown to be a necessary design decision.

\subsection{TabSyn as state-of-the-art reference}

TabSyn \cite{zhang2024tabsyn} is the top fidelity baseline, published at ICLR 2024. It is the method that RE-TabSyn builds directly upon. Including it is both an obligation and a deliberate positioning choice.

It is an obligation because any paper that proposes an extension of TabSyn must show its performance relative to the method it extends. Omitting TabSyn would be a major weakness in the comparison.

It is a positioning choice because TabSyn's performance (KS = 0.109, AUC utility ratio = 97.7\%) defines the ``fidelity ceiling'': the best possible distributional fidelity achievable with a VAE-plus-diffusion approach without any class conditioning overhead. RE-TabSyn accepts a modest fidelity cost (KS = 0.171) to gain class control. The TabSyn comparison makes the size of that cost visible and allows the reader to judge whether it is acceptable for their use case.

\subsection{What the four baselines together establish}

The four baselines are not redundant -- they each answer a different question:

\begin{enumerate}
    \item \textbf{CTGAN}: Is RE-TabSyn better than the current industry-standard GAN?
    \item \textbf{TVAE}: Is the VAE-plus-diffusion design better than a VAE alone?
    \item \textbf{TabDDPM}: Is latent-space diffusion better than direct diffusion?
    \item \textbf{TabSyn}: What is the fidelity cost of adding class control?
\end{enumerate}

Together they triangulate RE-TabSyn's position in the design space: it is better than CTGAN as the GAN representative, better than standalone TVAE, better than TabDDPM for direct diffusion, and adds class control that TabSyn, the current best-fidelity method, lacks.

TVAE's standalone results are not reported in every table in Chapter~5 because its performance is consistently dominated by both TabSyn and RE-TabSyn, and including it in every plot would add visual noise without adding information. Where TVAE results are relevant to a specific comparison (such as isolating the VAE family's performance), they are discussed explicitly.

\section{Identified Research Gaps}

The analysis of 151 publications presented in this chapter reveals three critical gaps that motivate RE-TabSyn.

\begin{table}[htbp]
\centering
\caption{Research Gap Analysis: Capabilities of Existing Methods vs.\ RE-TabSyn}
\label{tab:gap_analysis}
\small
\begin{tabularx}{\textwidth}{|l|C|C|C|C|C|}
\hline
\textbf{Capability} & \textbf{CTGAN} & \textbf{TVAE} & \textbf{TabDDPM} & \textbf{TabSyn} & \textbf{RE-TabSyn} \\
\hline
Mixed-type Support & Yes & Yes & Partial & Yes & Yes \\
\hline
High Fidelity (KS $<$ 0.20) & Yes & Yes & No & Yes & Yes \\
\hline
Stable Training & No & Yes & Yes & Yes & Yes \\
\hline
Mode Coverage & No & Partial & Yes & Yes & Yes \\
\hline
Class Control & No & No & No & No & \textbf{Yes} \\
\hline
Post-hoc Ratio Adjustment & No & No & No & No & \textbf{Yes} \\
\hline
Privacy Preservation & Weak & Weak & Moderate & Strong & Strong \\
\hline
\end{tabularx}
\end{table}

\begin{enumerate}
    \item \textbf{Gap 1: Absence of Controllable Minority Class Generation.} No existing method enables practitioners to specify desired minority ratios during generation. All methods faithfully replicate (or degrade) the training class distribution without offering any mechanism for adjustment. This is not a minor limitation -- it means that synthetic data cannot address class imbalance, which is the primary reason financial institutions would want to use it in the first place. A fraud detection team that generates synthetic data with TabSyn still ends up with the same 0.5\% fraud prevalence as their real data. They would need to generate 200 times as much data to get the same number of fraud cases as legitimate ones, which creates its own statistical problems.

    \item \textbf{Gap 2: CFG Application to Tabular Domain.} Classifier-Free Guidance has proven effective in image synthesis, becoming the default conditioning mechanism for leading text-to-image models. Its ability to control class distribution at inference time, without requiring retraining or auxiliary classifiers, makes it a natural fit for the minority ratio control problem. Yet its potential for structured tabular data class control remained entirely unexplored. Our systematic search across all 151 papers found zero instances of CFG applied to tabular synthesis. The adaptation is non-trivial: image CFG operates on continuous pixel values, while tabular data has mixed types, different scales, and complex inter-column dependencies that require careful architectural design to handle correctly.

    \item \textbf{Gap 3: Unified Fidelity-Privacy-Control Framework.} Existing works optimize for either fidelity (TabSyn), privacy (DP-CTGAN), or rare event handling (SMOTE-based methods). No single framework addresses all three dimensions simultaneously, leaving practitioners with an uncomfortable choice between competing objectives. TabSyn gives excellent fidelity but no privacy guarantees and no class control. DP-CTGAN gives formal privacy but poor fidelity because noise injection destroys distributional structure. SMOTE gives class balance but generates samples that may not be realistic. RE-TabSyn aims to occupy the intersection of fidelity, privacy, and control that previously had no occupant.
\end{enumerate}

RE-TabSyn directly addresses all three gaps by:
\begin{enumerate}
    \item Adopting latent space diffusion for robust mixed-type handling, borrowing TabSyn's core architectural insight.
    \item Introducing Classifier-Free Guidance to tabular synthesis for the first time, enabling controllable minority ratios through a single guidance scale parameter.
    \item Maintaining competitive fidelity and strong empirical privacy alongside controllability, providing a unified framework for the first time.
\end{enumerate}

The cost of these gains is a modest reduction in peak fidelity (KS = 0.171 versus TabSyn's 0.109), representing the trade-off between tight distributional matching and class-conditional control. Chapter~5 argues that this trade-off is acceptable, and in imbalanced settings favorable: the utility gain from balanced training data outweighs the small fidelity loss.

The theoretical foundations underpinning this work draw from VAE theory \cite{kingma2014vae}, copula methods \cite{nelsen2006copulas}, causal structure learning \cite{zheng2024causal}, pre-trained DP-VAEs \cite{jiang2023dp2vae}, federated synthesis \cite{rajotte2022fedsyn}, graph neural networks for tabular data \cite{zhu2023gnn}, optimal transport \cite{villani2009optimal, hallin2021ksample}, contrastive predictive coding \cite{oord2018representation}, statistical data privacy \cite{slavkovic2023statistical}, Renyi DP \cite{balle2020renyi}, privacy-preserving decision frameworks \cite{haney2024decision}, optimal privacy control \cite{xu2024optimal}, meta-learning for synthetic data \cite{triantafillou2022metalearning}, the chi-square test \cite{pearson1900criterion}, and imbalanced learning \cite{he2009imbalanced, chawla2002smote}. Recent advances include adversarial detector training \cite{liu2023adversarial}, deep generative model surveys \cite{jordon2022deepgen}, foundation models for tabular data \cite{van2024foundation, bommasani2022opportunities}, neural ODEs \cite{garcia2024neural}, relational structure learning \cite{liu2023goggle}, active learning from synthetic data \cite{haussmann2020imital}, quantum-classical hybrid models \cite{coyle2024quantum}, federated learning \cite{mcmahan2017fedavg}, and hierarchical text-conditional generation \cite{ramesh2022clip}. The evaluation and implementation rely on XGBoost \cite{chen2016xgboost}, t-SNE visualization \cite{maaten2008tsne}, the Transformer architecture \cite{vaswani2017attention}, Diffusion Transformers \cite{peebles2023dit}, layer normalization \cite{ba2016layernorm}, and dropout regularization \cite{srivastava2014dropout}.
